\documentclass[11pt,letterpaper]{article}
\usepackage[margin=1in]{geometry}
\usepackage{amsmath,amssymb,amsthm}
\usepackage{hyperref}
\usepackage{cleveref}
\usepackage{booktabs}

\numberwithin{equation}{section}
\renewcommand{\theequation}{S14.\arabic{equation}}

\title{Supplement S14: Quark Charges, Color, and Mass Ratios \\
\large Supporting Manuscript \S17}
\author{UDT Collaboration}
\date{}

\begin{document}
\maketitle

\begin{abstract}
We derive the quark electric charges, the number of colors, and the
inter-generation mass ratios from the angular structure on $S^2$.
The color multiplicity $N_c = (2\ell+1) = 3$ is the orbital degeneracy.
Quark charges follow from rank-2 parity blocks: $Q_d = -1/3$,
$Q_u = +2/3$.  Four inter-generation mass ratios are each $2\pi^2$
times a sector factor, all sub-$2\%$.  Verified by
\texttt{analysis/20\_quark\_masses.py}.
\end{abstract}

\tableofcontents

%----------------------------------------------------------------------
\section{Color as orbital orientation}
%----------------------------------------------------------------------

The Dirac equation on $S^2$ decomposes into angular channels labeled
by $\kappa$.  At each $|\kappa|$, the orbital angular momentum
$\ell = 1$ provides a 3-dimensional representation space with basis
states $|m_\ell\rangle$ for $m_\ell = -1, 0, +1$.  These three
orientations are identified with the three color states:
\begin{equation}
  N_c = 2\ell + 1 = 3\,.
\end{equation}
The $\mathfrak{su}(3)$ algebra that governs color interactions closes
exactly on this $\ell = 1$ subspace (Supplement~S16).

%----------------------------------------------------------------------
\section{Quark electric charges}
%----------------------------------------------------------------------

The rank-2 tensor operators on the $\ell = 1$ space have a parity
block structure.  The sign counting within each rank-2 block gives:
\begin{align}
  Q_d &= -\frac{1}{N_c} = -\frac{1}{3}\,,\\
  Q_u &= +\frac{2}{N_c} = +\frac{2}{3}\,.
\end{align}
The total $R$-ratio:
\begin{equation}
  R(u,d) = N_c(Q_u^2 + Q_d^2) = 3 \times \frac{5}{9} = \frac{5}{3}\,,
\end{equation}
matches the Standard Model exactly.

The charge assignment $Q_u = +2/3$ also follows from the bridge
identity: $\sqrt{B_1/B_3} = \sqrt{(1/18)/(1/8)} = \sqrt{4/9} = 2/3$.

%----------------------------------------------------------------------
\section{Inter-generation mass ratios}
%----------------------------------------------------------------------

The four ratios between adjacent generations follow a pattern:
each is $2\pi^2$ times a sector factor from the $j$--$\pi$--$\ell$
bridge transition.

\begin{center}
\begin{tabular}{lclrc}
\toprule
Ratio & Formula & Value & Expt. & Error \\
\midrule
$m_s/m_d$ & $2\pi^2$ & $19.74$ & $20.0$ & $-1.3\%$ \\
$m_c/m_u$ & $6\pi^4$ & $584.0$ & $587.5$ & $-0.6\%$ \\
$m_b/m_s$ & $9\pi^2/2$ & $44.41$ & $44.8$ & $-0.8\%$ \\
$m_t/m_c$ & $14\pi^2$ & $138.2$ & $136.0$ & $+1.6\%$ \\
\bottomrule
\end{tabular}
\end{center}

The common factor $2\pi^2$ is the baseline inter-generation step.
The sector factors $(1, 3\pi^2, 9/4, 7)$ each involve the
multiplicities $(2, 3, 5, 7)$:
\begin{itemize}
  \item $m_s/m_d$: baseline, no sector factor.
  \item $m_c/m_u = 6\pi^4 = (2\ell+1)(2j+1) \times (2\pi^2)^2$:
    spin-orbital multiplied by squared step.
  \item $m_b/m_s = (9/4)\pi^2 = [(2\ell+1)/(2j+1)]^2 \times \pi^2$:
    squared orbital/spin ratio.
  \item $m_t/m_c = 14\pi^2 = 2(2|\kappa_\mathrm{max}|+1) \times \pi^2$:
    twice the orbit size.
\end{itemize}

\subsection{Base quark masses}

The lightest quarks have masses set by the selection rule for rank-2
bilinears --- squared multiplicity:
\begin{align}
  m_u &= (2j+1)^2\,m_e = 4\,m_e = 2.04~\text{MeV}
    \qquad(-5.4\%)\,,\\
  m_d &= (2\ell+1)^2\,m_e = 9\,m_e = 4.60~\text{MeV}
    \qquad(-1.5\%)\,.
\end{align}
These are approximate because $\overline{\text{MS}}$ scheme masses at
$\mu = 2$~GeV carry renormalization scheme dependence at the $\sim5\%$
level for the lightest quarks.

%----------------------------------------------------------------------
\section{Verification}
%----------------------------------------------------------------------

All results verified by \texttt{analysis/20\_quark\_masses.py}.
Output: \texttt{data/generated/20\_quark\_masses.json}.

\end{document}
